%======================================================================
\section{Protocol}\label{sec:protocol}
%======================================================================

The protocol rests on one rule. Uniformly collected data are read only to choose among trained models
and to measure them. They never enter a training objective, a nuisance model that a trained model uses,
or the stopping of training. This section states the controls that enforce the rule, how the splits,
the selection and the seeds are built, how the tests and sensitivity margins are defined, and how the
stored results were reconciled before the analysis. The methods to which the protocol is applied are
listed in Section~\ref{sec:methods}.

\subsection{Where each data source may enter}\label{sec:leakage}

Table~\ref{tab:leakage} lists the stages of the pipeline and the data each stage may read. ``Log'' is the
biased training log, which holds the exposure matrix and the outcomes of the exposed pairs.
``Validation'' and ``test'' are the parts of the uniformly collected data that a split assigns to each role.

\begin{table}[htbp]\centering\footnotesize
\caption{Permitted uses of each data source. The Coat row for the propensity is the one exception to the
rule, and it is discussed in the text.}\label{tab:leakage}
\begin{tabular}{lccc}
\toprule
Stage & Log & Validation & Test \\
\midrule
Exposure propensity $\hat p$ (Coat: shipped, see text) & yes & no & no \\
Additive imputation $\hat Y$, prior $z$, weights $Q$ & yes & no & no \\
Cross-fitted nuisances of the DRGP prior & yes & no & no \\
Training of every model & yes & no & no \\
Choice of configuration and epoch within a split & no & yes & no \\
Early stopping within a split & no & yes & no \\
Reported test score of the chosen model & no & no & yes \\
\bottomrule
\end{tabular}
\end{table}

The additive imputation is fitted on the log without cross-fitting, so the same log supplies $\hat Y$ and the
training data. The DRGP prior of \ref{sec:drgp} is cross-fitted over ten folds before its
coefficients are fitted. Test scores are computed after every epoch and stored, but no choice reads them.
The sensitivity analysis of the results reads them only to measure what selection on test data would
change, and it does not report such a selection as a method.

\paragraph{The Coat exception.}
The shipped Coat propensities were estimated by the dataset authors from a uniformly collected
sample~\citep{schnabel2016}. That sample is also the source of our validation and test data, so uniform
data reach a training input on Coat, indirectly. Every propensity-based row on Coat uses these
propensities: DR-LightGCN, NAVIP, DR-JL, MRDR, DR adjacency, DRUP and the prior rows. To check the effect
on the prior rows, we refitted them with propensities estimated from the log by the logistic model of
Section~\ref{sec:prior}. Table~\ref{tab:T10} gives the results. Relative to their no-prior backbones, the
prior gains are $+19.44\times10^{-3}$ [18.66] for iALS-p and $+33.62\times10^{-3}$ [16.72] for LightGCN-p
with estimated propensities. With the shipped propensities they are $+16.82\times10^{-3}$ [18.04] and
$+31.98\times10^{-3}$ [20.65]. The paired differences between the two propensity inputs are
$+2.63\times10^{-3}$ [5.27] for iALS-p and $+1.64\times10^{-3}$ [10.30] for LightGCN-p, with
$p_\mathrm{Holm}=1.000$ in both cases. The values in brackets are 95\% half-widths. The stored
no-prior runs of iALS and LightGCN serve as the reference for both settings, because neither reads the
propensities in its loss. The check changes only the propensity input of the prior rows. For LightGCN-p the
gain with estimated propensities is not removed. For iALS-p the estimated-propensity gain is not significant
($p_\mathrm{Holm}=0.33$), so the statement that the gain remains is a point estimate only. The check does not
test the assumption that exposure is independent of the outcome given the covariates, which the closed-form
posterior of the method requires. It also does not cover the other propensity-based rows of Coat: DR-LightGCN,
NAVIP, DR-JL, MRDR, DR, DRUP and the QP-LightGCN family were not refitted, so their dependence on the shipped
propensities is not checked.

\subsection{Splits}\label{sec:splits}

On Coat, the split is made within each user: 30\% of each user's candidate items go to validation and the
rest go to test. On Yahoo!\,R3, KuaiRand-Pure and KuaiRec, the split is made over users: 30\% of the test
users go to validation, and no user is in both parts of a split. Each split is a function of its seed.
Coat uses ten splits (seeds 0 to 9) and the other datasets use five (seeds 0 to 4). On Yahoo!\,R3, a
fixed random 5\% of the users (270) is excluded from validation and testing.

A user is a test user of a split only if the user has a test entry in that split. A user can fall in the
validation part of every split, so the number of test users in the files of the final method is 225,
2{,}310, 8{,}321 and 1{,}410 on Coat, Yahoo!\,R3, KuaiRand-Pure and KuaiRec. These are slightly below the
evaluated-user counts of Table~\ref{tab:datasets}, and equal to them on Coat.

Candidates are ranked for each user within the part being scored, and users with no positive candidate in
that part are skipped. We rank all candidates and do not sample negatives, because sampled metrics do not
agree with the full ranking metric~\citep{krichene2020sampled}. The metric is nDCG@5 on Coat and
Yahoo!\,R3, nDCG@10 on KuaiRand and nDCG@20 on KuaiRec. Tables report the mean over splits and the standard
deviation over splits with $\mathrm{ddof}=0$. The ensemble stage of \ref{sec:ensemble} asserts, for
every split it reads, that validation and test users are disjoint (on Yahoo!\,R3, KuaiRand and KuaiRec) or
that validation and test items are disjoint for each user (on Coat).

\subsection{Training and per-split selection}\label{sec:selection}

Each configuration is trained once on the log, and the same trained model is scored on every split. The
splits therefore share training and differ only in selection. After each epoch, the model is scored on the
validation and test parts of every split. Each split keeps the (configuration, epoch) pair with the highest
validation nDCG, and the first maximum wins a tie. A split stops when its validation score has not improved
for the patience period: five epochs on Coat and Yahoo!\,R3, and three epochs on KuaiRand and KuaiRec. The
shared run ends when every split has stopped or the epoch limit is reached. Once a split has stopped, its
selection is fixed. Later epochs are recorded but cannot be chosen.

The gradient-trained models use epoch limits of 100 (Coat), 60 (Yahoo!\,R3), 40 (KuaiRand) and 30 (KuaiRec).
Batch sizes are 256, 1{,}024, 8{,}192 and 32{,}768. Within a dataset, all methods share the epoch limit,
batch size and patience, except AdvInfoNCE on KuaiRec, which uses batch size 4{,}096. On KuaiRec, each
training epoch is capped at 2{,}000{,}000 pairs.

Three further quantities are computed from the same per-epoch records for the sensitivity analysis. Each is defined on the
rows that the stopping rule of its split allows to be chosen.
\begin{itemize}
\item The per-split test oracle is the largest test value of a split over all configurations and epochs.
It chooses on test data.
\item The configuration-level oracle is the configuration whose mean test value, at its own
validation-selected epoch, is largest. It also chooses on test data.
\item The expected maximum of the validation-selected test value over $k$ configurations is computed as a
function of $k$~\citep{dodge2019showyourwork}, and the share of 1{,}000 random subsets of size $k_{\min}$ in
which the sign of a primary comparison flips is reported. Here $k_{\min}$ is the smaller of the two configuration
counts of the comparison.
\end{itemize}
The first two are diagnostics. They are never reported as the result of a method.

The reconciliation of the stored results, which checks that every reported selection can be recomputed from its per-epoch records, is described in \ref{sec:reconcile}.

\subsection{Seeds}\label{sec:seeds}

Each split's selected configuration is retrained with seeds 1, 2 and 3. Each split re-selects its epoch on
its own validation part, and the configuration is not changed. The retrained final method is compared with
the two baselines that had the highest main-run test means in each dataset: iALS and DR-LightGCN on Coat,
r-AdjNorm and LightGCN on Yahoo!\,R3 and KuaiRand, and DR-JL and MRDR on KuaiRec. The baselines are retrained
in the same way. Because these two baselines were chosen on test data, the seed comparisons are descriptive,
and they are reported without a significance claim beyond the Holm-adjusted values given in the text; the
Holm correction runs over the two tests of each dataset. DRUP is deterministic and is not retrained.
The seed-0 values in the tables are the main-run values.

Seeds 1 to 3 were not run for the Coat prior rows (iALS-p and LightGCN-p). Their Coat values come from
seed 0 only. The result files also contain re-runs of split 0 with seeds 100 to 102, which are stored
separately. These runs are not used in any table.

\subsection{Statistical tests}\label{sec:stats}

Differences are computed per user. For each user, the test score is averaged over the splits in which the
user is a test user. For a comparison $A-B$, let $d_u$ be the per-user difference for user $u$. The mean
$\bar d$ over users is the reported difference, and a paired $t$ test over users gives the two-sided $p$-value
and a 95\% half-width. Differences are reported as $\times10^{-3}$ with the half-width in brackets. The
reported difference is not the difference of the two table means. The two can have opposite signs. For
example, on Coat the difference of the table means of DR-LightGCN and LightGCN is $+4.1\times10^{-3}$, while
the mean per-user difference is $-0.65\times10^{-3}$ [30.16]. The per-user difference is the statistic used
throughout.

A difference is called significant only when its Holm-adjusted $p$ is below 0.05. The families are as follows.
\begin{description}
\item[Primary family, per dataset.] Four comparisons, P1 to P4, are defined in code from the stored
results. P1 and P4 choose their third-party debiaser and trained comparator, respectively, by mean
validation score, and never by test score.
  \begin{itemize}
  \item P1: the third-party debiaser with the highest mean validation score on the dataset, against LightGCN.
  \item P2: each simple prior against its backbone, that is iALS-p against iALS and LightGCN-p against
  LightGCN, where both exist on the dataset.
  \item P3: the best refinement against the simple prior of the same learner. On Coat this gives two
  comparisons, the QP-LightGCN family against LightGCN-p and DR-iALS against iALS-p. On KuaiRec it gives the
  QP-LightGCN family against LightGCN-p.
  \item P4: DRUP against the trained model with the highest mean validation score. This is computed on
  KuaiRec only.
  \end{itemize}
  The primary families contain five comparisons on Coat and KuaiRec and two on Yahoo!\,R3 and KuaiRand.
\item[Walk-correction family.] DRUP against DR adjacency, one comparison per dataset.
\item[Refinement families.] Up to three families per dataset. The first holds each refinement against the simple
prior of the same learner. The second, on Coat only, holds DR-iALS against iALS-p. The third holds the shuffled
and popularity controls against LightGCN, on the datasets where those controls were run.
\item[Backbone family.] Every method against LightGCN, reported in Table~\ref{tab:T12-coat} to
Table~\ref{tab:T12-kuairec} in the appendix.
\end{description}

The Nadeau and Bengio correction~\citep{nadeau2003} is applied as a second test to the split means. With $J$ splits
and split-level differences $e_j$, the variance is $\hat\sigma^2 = s_e^2\,(1/J+n_2/n_1)$, where $s_e^2$ is the sample
variance of the $e_j$. The statistic $\bar e/\hat\sigma$ is compared with a $t$ distribution with $J-1$ degrees of
freedom, and Holm adjustment is applied to these second $p$-values within the same families. In the split design,
the selection part is $n_1$ and holds 30\% of the users (or of each user's candidates), and the test part is $n_2$
and holds the remaining 70\%. The code uses $n_2/n_1=(1-\rho)/\rho=0.7/0.3=7/3$ with $\rho=0.3$, which is the
standard ratio with the selection part in the role of the training part. The standard derivation does not apply
directly, because the models are trained once on the log and only the selection and test parts vary. With $J=5$ the
variance factor is $1/5+7/3\approx2.53$. Reading the same 70/30 proportions the other way (test as $n_1$) would give
$n_2/n_1=3/7$ and a factor of $1/5+3/7\approx0.63$. The larger factor is used, so the correction is conservative
relative to that reading. For Coat with $J=10$ the factor is $1/10+7/3\approx2.43$ (against $0.53$). We report the
corrected values beside the standard paired test and do not treat either as exact.

Equivalence is assessed with two one-sided tests~\citep{schuirmann1987} at margin $\delta$, each at level 0.05. The
two null hypotheses are $\bar d\le-\delta$ and $\bar d\ge\delta$. The reported TOST $p$-value is the larger of the two
one-sided $p$-values. The margins are $\pm0.010$ nDCG@5 on Coat and Yahoo!\,R3, and $\pm0.003$ nDCG on KuaiRand
(nDCG@10) and KuaiRec (nDCG@20). They are sensitivity thresholds, chosen before the analysis and not derived from
the literature, and TOST is presented as a sensitivity analysis of these thresholds. On Yahoo!\,R3, KuaiRand and
KuaiRec, a TOST result below 0.05 is stated as equivalence within the margin. On Coat, the same result is stated
only as the exclusion of differences larger than the margin in absolute value, and equivalence is not claimed. A
result inside a margin is never stated as equality.

The minimum detectable effect at 80\% power and two-sided level 0.05 is
$\mathrm{MDE}=(t_{0.975,\,n-1}+t_{0.80,\,n-1})\,s_d/\sqrt{n}$, where $n$ is the number of users, $s_d$ is the
standard deviation of the per-user differences, and $t_{q,\nu}$ is the $q$ quantile of a $t$ distribution with
$\nu$ degrees of freedom. The MDE uses only the user-level standard deviation and ignores the dependence between
splits, so it is optimistic. For the primary comparisons in the analysis, the MDE is $11.8$ to $43.1\times10^{-3}$
on Coat, $4.7$ to $5.6\times10^{-3}$ on Yahoo!\,R3, $1.5$ to $1.9\times10^{-3}$ on KuaiRand and $0.7$ to
$4.9\times10^{-3}$ on KuaiRec. The Coat values are large because the Coat analysis has 225 evaluated users,
and an MDE of this size bounds what the Coat data can show.

Each user contributes one difference per comparison, which is the mean over the splits in which the user is a
test user. The splits share users and candidates, so the paired $t$ test over overlapping splits can be
anti-conservative~\citep{nadeau2003}. For this reason the Nadeau and Bengio column is reported beside it for the
primary comparisons. The tests treat the training log, the fitted nuisances and the trained models as fixed.
